We study a growth model for causal sets in which every new element is
inserted into the interior of the existing causal set, so that the
density of elements increases dynamically. The insertion rule is Lorentz
invariant: a child is placed at a small, Poisson-distributed proper-time
separation from a uniformly chosen parent, with a rapidity drawn from
the Haar measure \(\sinh^2\eta\,d\eta\), optionally damped by a
Maxwell--Jüttner factor \(e^{-\alpha\cosh\eta}\) whose rest frame is
read off the state. We use the model to address four open problems of
causal set theory.

\emph{Continuum limit.} The empirical measure of the growing set
converges almost surely for every \(\alpha\ge 0\), and every finite
induced-subposet density converges, so the process has a poset limit: a
uniform sprinkling of a conformally flat spacetime. For \(\alpha=0\) the
limit is explicit, \(a^4\propto t^2+(T-t)^2\).

\emph{Mesoscopic manifoldlikeness.} Exact second-moment recursions, an
eigenvalue criterion and martingale concentration give high-probability
versions of Madsen's conditions F2 (counting equals volume, uniformly
over all mesoscopic causal diamonds) and F3 (longest chain equals proper
time) for the conformal metric. Density clumping, a Pólya-urn effect, is
confined to wavelengths above \(8.0\)--\(9.1\) damping lengths.

\emph{Lorentz invariance versus locality.} Frame independence holds at
the 1--5\% resolution of intrinsic tests. Lorentz violation enters only
through the clumping excess and is negligible in the physical regime.

\emph{Rogue posets and quantum dynamics.} Order-only insertion rules
fail in three distinct ways, all traced to the light-cone nonlocality of
order-space neighbours, whose symmetric difference grows as \(N^{1/2}\)
in 2D and \(N^{0.8}\) in 4D. Attaching the phase
\(e^{i\beta S_{\rm BD}}\) to such histories, we prove a defect-gas
theorem: for \(M\) independent local defect sites with \(K\)
alternatives and uniformly distributed relative phases, the defect
amplitude grows as \(e^{M\,m(1+x_1+\dots+x_K)}\), where \(m\) is the
Mahler measure. The rate is zero for \(K=1\) and positive for
\(K\ge 2\), so multi-type local defects are not suppressed by
interference. Numerical checks support the uniform-phase hypothesis at
generic \(\beta\), additivity, the overlap correction and the absence of
double counting. The key analytic steps are machine-checked in Lean 4
with Mathlib.
